\documentclass[a4paper,12pt]{article}
\usepackage{geometry}
\geometry{left=2.5cm,right=2.5cm,top=2.5cm,bottom=2.5cm}
\renewcommand{\textfraction}{0.15}
\renewcommand{\topfraction}{0.85}
\renewcommand{\bottomfraction}{0.65}
\renewcommand{\floatpagefraction}{0.60}
\usepackage{amsmath}
\usepackage{amsfonts}
\usepackage{mathrsfs}
\usepackage{amsthm}
\usepackage{extarrows}
\usepackage{bm}
% \newcommand{\bm}{\symbfit}    % `bm` confilicts with `unicode-math`. In that case use \symbfit for bold math symbols
\usepackage{graphicx}
\usepackage[section]{placeins}
\usepackage{flafter}
\usepackage{array}
\usepackage{caption}
\usepackage{subcaption}
\usepackage{color}
\usepackage{multirow}
\usepackage{natbib}
% \usepackage{enumerate}
\usepackage{enumitem}    % more flexible than `enumerate` package, the reference will carry the whole label appearance, not just the counter, unlike the `enumerate` package.
\usepackage{upgreek}    % 'upgreek' letters

% \pdfstringdefDisableCommands{\let\bm=\relax}

\newtheorem{thm}{Theorem}
\newtheorem{cor}{Corollary}
\newtheorem{assum}{Assumption}
\newtheorem{rem}{Remark}
\newtheorem{lem}{Lemma}
\newtheorem{prop}{Proposition}

\setcounter{topnumber}{5}    % Maximum number of floats that can appear at the top of a text page; default 2. 
\setcounter{bottomnumber}{5}   % Maximum number of floats that can appear at the bottom of a text page; default 1. 
\setcounter{totalnumber}{10}    % Maximum number of floats that can appear on a text page; default 3. 

\DeclareMathOperator*{\argmaxdown}{arg\,max}
\DeclareMathOperator*{\argmindown}{arg\,min}
\DeclareMathOperator{\argmax}{arg\,max}
\DeclareMathOperator{\argmin}{arg\,min}

% cross-reference to other files
% the externaldocument should be compiled 
% (and at least twice if you're using xr-hyper)
% \usepackage{xr-hyper}
% \usepackage{xr}

% --- external document (ordinary setting) ---
% \externaldocument{external_tex_file}
% --- end of ordinary setting ---

% --- external document (overleaf setting) ---
% externaldocument settings for Overleaf
% \makeatletter
% \newcommand*{\addFileDependency}[1]{% argument=file name and extension
% \typeout{(#1)}
% \@addtofilelist{#1}
% \IfFileExists{#1}{}{\typeout{No file #1.}}
% }
%   \makeatother
%   \newcommand*{\myexternaldocument}[1]{%
%   \externaldocument{#1}%
%   \addFileDependency{#1.tex}%
%   \addFileDependency{#1.aux}%
% }
%   \myexternaldocument{external_tex_file}
%   --- end of overleaf setting ---

%   mathtools can be used to define labeling format for equations
%   one can use \eqref for a reference to a labeled equation.
\usepackage{mathtools}
\newtagform{supp}{(S-}{)}    % define a equation labeling format for suppliment
\usetagform{supp}            % use the supp format
\usetagform{default}         % use the default format

\usepackage{algorithmic}
\usepackage{algorithm}
\renewcommand{\algorithmicrequire}{\textbf{Input:}}
\renewcommand{\algorithmicensure}{\textbf{Output:}}

% package for hyperlinks
% It's error-prone because hyper link is quite difficult
% due to the fact the typesetting environment is complex.
% So you can disable this package and finish the document.
% Then sort out the hyperlink thing.
\usepackage[colorlinks,linkcolor=red,anchorcolor=blue,citecolor=green,CJKbookmarks=True]{hyperref}

% package for displaying highlighted codes
\usepackage{minted}

% package for input codec and output rendering font
\usepackage[utf8]{inputenc}    % it is always good practice to use utf8
% you can also try latin1, latin2, cp1252 and cp1250
\usepackage[T1]{fontenc}    % the default is T0, which only contains 128 characters
% you can try T1, T2A, T2B
% you can refer to https://www.overleaf.com/learn/latex/international_language_support#Font_encoding
% for more details.


\title{Only One Event in Cox Model}
\author{Chao Cheng}
\date{\today}



\begin{document}
\maketitle
\tableofcontents{}

\section{Introduction}
\label{sec:introduction}

Let's consider the Cox model and there are \textbf{No tied} observations. And define the distinct times of failure (or censoring) $\tau_1 < \tau_2 < \cdots < \tau_n$. Denote
\[
  R_j = \left\{i\middle|U_i \geq \tau_j\right\} = \text{risk set at $\tau_j$}, 
\]
where $U_i = \mathbf{min}\left(T_i, C_i\right)$ is the observed time and
\[
  Z_{(j)} = \text{value of $Z$ for the subject who {\color{red}fails} at $\tau_j$}. 
\]
The partial likelihood $L_1(\beta)$ is defined as
\begin{equation}
  \label{eq:partial_likelihood_l1}
  L_1\left(\beta\right)
  \overset{\Delta}{=}
  \prod_j\left\{
    \frac{e^{\beta Z_{{\color{red}\left(j\right)}}}}{
      \sum\limits_{l\in R_{j}}e^{\beta Z_{{\color{red}l}}}
    }
  \right\}
\end{equation}
In this notes, we will consider the clinical trial scenario. So $Z_i\in\left\{0, 1\right\}$ for $i = 1, \cdots, n$. $Z_i = 0$ denotes control group and $Z_i = 1$ denotes the treatment group. And we will consider the case when there's only one event observed.

\section{Only One Event Observed}
\label{sec:only-one-event}

If there's only one event observed, then \eqref{eq:partial_likelihood_l1} can be simplified as
\begin{equation}
  \label{eq:partial_likelihood_only_one_event}
  L_1\left(\beta\right)
  = \frac{e^{\beta Z_j}}{\sum\limits_{l\in R_j}e^{\beta Z_l}}
  = \frac{e^{\beta Z_j}}{n_{trt}e^{\beta\times 1} + n_{ctrl}e^{\beta \times 0}}
  = \frac{e^{\beta Z_j}}{n_{trt}e^{\beta} + n_{ctrl}}, 
\end{equation}
where $n_{trt}$ is the number of observations in $R_j$ that belongs to treatment group and $n_{ctrl}$ is the number of observations in $R_j$ that belongs to the control group. 
\par
To maximize \eqref{eq:partial_likelihood_only_one_event}, we can see
\[
  L_1\left(\beta\right)
  = \left\{
    \begin{aligned}
      & \frac{e^{\beta}}{n_{trt}e^{\beta} + n_{ctrl}}
        = \frac{1}{n_{trt}}
        - \frac{n_{ctrl}}{n_{trt}}\frac{1}{n_{trt}e^{\beta} + n_{ctrl}}
      && \;\text{if } Z_j = 1 \\
      & \frac{1}{n_{trt}e^{\beta} + n_{ctrl}}
      && \;\text{if } Z_j = 0
    \end{aligned}
  \right.
\]
If $Z_j = 1$, \eqref{eq:partial_likelihood_only_one_event} is maximized at $\beta = \infty$. If $Z_j = 0$, \eqref{eq:partial_likelihood_only_one_event} is maximized at $\beta = -\infty$. 




\bibliographystyle{plainnat}
\bibliography{../../ref}





\end{document}


%%% Local Variables:
%%% mode: latex
%%% TeX-master: t
%%% End:
